% =====================================================================
%  範例文件：莫蘭迪 × ADHD 友善閱讀模板
%  編譯：latexmk -xelatex main.tex
%
%  可用選項（加在 \documentclass[...] 裡）：
%    dark      深色護眼模式
%    print     列印模式（白底）
%    compact   較緊湊的行距
%    nobionic  關閉英文閱讀引導粗體
% =====================================================================
\documentclass{morandi-adhd}

\title{莫蘭迪 ADHD 友善筆記模板}
\subtitle{低刺激配色、清楚分塊、一眼找到重點}
\author{你的名字}
\date{\today}
\readingtime{約 5 分鐘}

\begin{document}

\maketitle

\begin{tldr}
  這份模板用\key{柔和的莫蘭迪色}降低視覺疲勞，用\key{固定的顏色語意}和\key{短段落}幫助專注。
\end{tldr}

\tableofcontents

\section{為什麼需要 ADHD 友善排版？}

\begin{goals}
  \begin{checklist}
    \item 了解這份模板的設計原則
    \item 知道每一種顏色方塊代表什麼
    \item 能開始用它寫自己的筆記
  \end{checklist}
\end{goals}

注意力不容易維持的讀者，面對\hl{一大片密密麻麻的文字}時，常常讀了三行就開始分心。
這份模板把文件拆成\key{一小塊一小塊}，每一塊都有明確的顏色和標題。

\begin{keypoint}
  一段只講一件事。段落短、留白多，讀者就比較容易「讀完一塊、休息、再讀下一塊」。
\end{keypoint}

\subsection{設計原則}

\begin{itemize}
  \item \key{低彩度配色}：莫蘭迪色帶有灰調，不會搶走注意力。
  \item \key{暖白背景、深灰文字}：比純白底純黑字更不刺眼。
  \item \key{加大行距、靠左對齊}：避免字距忽大忽小，視線不易跳行。
  \item \key{閱讀進度條}：頁尾的綠色條告訴你「還剩多少」。
    \begin{itemize}
      \item 知道終點在哪裡，比較不會焦慮。
    \end{itemize}
\end{itemize}

\takeabreak

\section{顏色語意一覽}

每種方塊都有固定的顏色，看久了大腦會自動辨認，不需要每次重新判讀。

\begin{definition}[莫蘭迪色]
  源自義大利畫家 Giorgio Morandi 的靜物畫，特色是在顏色中加入灰色，形成低飽和、溫和的色調。
\end{definition}

\begin{example}
  \begin{steps}
    \item 打開 \texttt{main.tex}
    \item 把標題、作者改成你自己的
    \item 用 \texttt{latexmk -xelatex main.tex} 編譯
  \end{steps}
\end{example}

\begin{tip}
  重要的詞用 \verb|\key{...}|，想畫螢光筆用 \verb|\hl{...}|，次要資訊用 \verb|\soft{...}| \soft{（像這樣變淡）}。
\end{tip}

\begin{warning}
  這個模板需要用 \key{XeLaTeX} 編譯，pdfLaTeX 無法處理中文字型。
\end{warning}

\begin{question}
  你讀完這一節，還記得「綠色邊線的方塊」代表什麼嗎？
\end{question}

\subsection{講義常用方塊}

\begin{bigidea}
  整章只記一件事的話，就記這個框裡的內容。
\end{bigidea}

\begin{terms}
  \key{莫蘭迪色}：加入灰色、飽和度低的顏色。\\
  \key{色盤}：一組固定搭配使用的顏色，定義在 \texttt{morandi.sty}。
\end{terms}

\begin{derivation}[面積公式]
  長方形面積 $= \text{長} \times \text{寬}$，三角形是同底同高長方形的一半，所以
  $A = \tfrac{1}{2}bh$。
\end{derivation}

\begin{selfcheck}
  \begin{checklist}
    \item 我能說出重點框和定義框的差別
    \item 我知道答案字母用 \verb|\answer{B}| 標示，例如：答案 \answer{B}
  \end{checklist}
\end{selfcheck}

\subsection{表格}

\begin{center}
  \zebra
  \begin{tabular}{lll}
    \toprule
    \headerrow \headcell{環境} & \headcell{顏色} & \headcell{用途} \\
    \midrule
    \texttt{keypoint}   & 鼠尾草綠 mdSage & 重點     \\
    \texttt{definition} & 乾燥玫瑰 mdTerm & 定義     \\
    \texttt{example}    & 灰藍 mdBlue     & 例題     \\
    \texttt{goals}      & 霧藍灰 mdDeep   & 本節目標 \\
    \texttt{tip}        & 奶茶 mdSand     & 小提示   \\
    \texttt{warning}    & 淺玫瑰 mdRose   & 注意     \\
    \texttt{question}   & 暖灰 mdMuted    & 想一想   \\
    \texttt{terms}   & 奶茶 mdSand     & 名詞     \\
    \bottomrule
  \end{tabular}
\end{center}

\chunkbreak

\subsection{英文閱讀引導}

英文段落可以用 \verb|\bionic{...}|，每三個單字中，第一個單字的前兩個字母會加粗，當作視線的「落腳點」：

\bionic{Reading long paragraphs can be exhausting when attention keeps drifting away from the page.}

\takeabreak[做得好！只剩最後一節]

\section{數學也可以}

行內公式 $E = mc^2$，或是獨立公式：
\[
  \int_0^1 x^2 \, dx = \frac{1}{3}
\]

\begin{summary}
  \begin{checklist}
    \done 用顏色區分資訊類型
    \done 短段落 + 留白
    \done 進度條減少焦慮
    \item 開始寫你自己的筆記！
  \end{checklist}
\end{summary}

\end{document}
